\SourceLeaf{359}%
geodesic line of known length and direction has long been treated in many ways. In choosing a method suited to one's needs, one must indeed consider various circumstances. For example, it matters whether the problem has to be solved for just one or a few specific cases or for very many. In the latter case it is important that the method make each final calculation as convenient and easy to follow as possible, even if applying the method may first require some general preparatory work. Equally important is whether all the results of an extensive trigonometrical survey are wanted in the form of geographical longitude and latitude, and \textit{only} in \textit{that} form, or whether the positions of all the points are also given in another form, for example in rectangular coordinates. In the latter case there is less need to give the geographical position with the very highest accuracy.

The formulas given by \textsc{Dusejour}, \textsc{Legendre}, \textsc{Delambre} and others take account of only the first power of the flattening. This will scarcely matter in practice, since the flattening of the Earth's spheroid is after all only a small fraction. The new method works independently of how small the flattening is. The intention here is to attach no great practical importance to that advantage. The better of those methods may well give sufficient accuracy in most cases, although an adequate mathematical demonstration is lacking. It can be said, however, that when the necessary prerequisites are at hand, the new method gives a more convenient calculation whose essential content can be set out in considerably less space. \textsc{Bessel}'s solution of 1825 bears the mark of great mathematical refinement and is entirely independent of the assumption that the distance between the two points is small compared with the dimensions of the whole terrestrial spheroid. This is undoubtedly an advantage of the method from a theoretical point of view. In judging its practical value, however, the following circumstances must be considered. The method makes no distinction at all between large and smaller distances. It requires calculations of the same length in every case,
